% !TeX root = ../Masterarbeit.tex

\documentclass[Masterarbeit.tex]{subfile}

\begin{document}
	
	\chapter*{Abstract}
	
	Slowly varying ocean heat content is one of the most important variables when describing the Earth system's memory on interannual to decadal time scales. Numerical climate models simulate the evolution of oceanic heat content after incorporating a set of observations via model assimilation. However, incomplete observations, particularly in the subsurface ocean constrain the accuracy of model assimilation and lead to large uncertainties. Kadow et al (2020) have proven that artificial intelligence can successfully be utilized to reconstruct missing climate information in the form of surface temperatures. In the following, I investigate the possibility to train their neural network to reconstruct missing ocean heat content data. \\
	To perform the reconstruction, I use a three-dimensional convolutional neural network originally designed for image inpainting technology. Combining artificial intelligence with earth system modeling, the network is trained and tested with a 16 member ensemble model assimilation from the Max-Planck Institute Earth System Model (MPI-ESM). These model runs are masked mirroring real life observations and can be adapted in their missing value density to the reconstruction time period in question. The monthly training and testing datasets, which are fed individually to neural network, run from January 1958 till December 2020 and thus incorporate an additional time dimension into the training process. \\
	
	
\end{document}